\section{Related Work}
\label{submission:sec:related}

\subsection{Model Merging and Weight-Space Consolidation}
Model merging has emerged as a powerful paradigm for multi-task learning without the computational overhead of joint training or multi-task fine-tuning \cite{wortsman2022model, choshen2022fusing}. The foundational concept relies on the empirical observation that neural networks fine-tuned from the same pre-trained initialization lie within a shared low-loss basin, permitting linear interpolation of their weights \cite{ainsworth2022git, wortsman2022model}. Ilharco et al. \cite{ilharco2022editing} formalized this via \emph{Task Arithmetic}, showing that adding task-specific "task vectors" to a pre-trained backbone creates multi-task merged models. Subsequent works have introduced heuristic filtering to mitigate parameter interference (e.g., TIES-Merging \cite{yadav2023ties} and ZipIt! \cite{stoica2023zipit}) or learning layer-wise blending coefficients via test-time gradient-based optimization on small validation subsets (e.g., AdaMerging \cite{yang2023adamerging}). Despite these developments, traditional model-merging literature operates under the implicit assumption of high-precision (FP16 or FP32) representations, ignoring the physical memory and hardware latency constraints that dictate real-world deployment.

\subsection{Post-Training Quantization (PTQ)}
To facilitate the deployment of deep neural networks on resource-constrained edge devices, neural network quantization compresses continuous floating-point weights into discrete low-bit integer formats (typically INT8 or INT4) \cite{nagel2021whitepaper, gholami2021survey}. Post-Training Quantization (PTQ) is highly favored in industrial settings as it eliminates the need for expensive quantization-aware training (QAT) on complete datasets \cite{krishnamoorthi2018quantizing}. Modern PTQ techniques focus on minimizing the discretization error through adaptive rounding mechanisms (AdaRound \cite{nagel2020adaround}), outlier-aware scale adjustments (SmoothQuant \cite{xiao2023smoothquant}, AWQ \cite{lin2023awq}), or block-wise reconstruction objectives (BRECQ \cite{li2021brecq}). While PTQ is highly effective for single-task networks, applying PTQ to a multi-task merged model is exceptionally challenging. The consolidated parameter space is highly non-linear, and the quantization operator introduces non-smooth discretization noise that destroys the fragile weight-space alignment required to preserve joint expert performance.

\subsection{Quantization-Aware Model Merging}
Recognizing the limitations of naive post-hoc quantization, recent pioneering works have attempted to bridge the gap between weight fusion and low-bit quantization. The standard baseline, Q-Merge \cite{qmerge2025}, optimizes layer-specific merging coefficients $\Lambda$ directly under a simulated target quantization schema. Because the rounding function has zero gradients almost everywhere, Q-Merge employs the Straight-Through Estimator (STE) \cite{bengio2013ste, yin2019understanding} to propagate gradients backward and update $\Lambda$ using first-order optimizers like Adam \cite{kingma2014adam}. The optimization is driven by unsupervised objectives, primarily prediction entropy minimization \cite{wang2020tent, zhang2021memo}, evaluated on a small, local test-time calibration stream. Sibling methods such as PolyMerge \cite{polymerge2025} parameterize layer coefficients using low-degree polynomials to constrain the optimization space, while RegCalMerge \cite{regcalmerge2025} proposes Class-Capacity Normalization and Elastic Spatial Regularization to stabilize the adaptation process.

\subsection{Our Critical Departure: Methodological Skepticism}
While these advanced quantization-aware merging frameworks claim outstanding low-bit performance, their evaluation protocols suffer from severe methodological limitations. First, they evaluate their methods under the exact same quantization schema used during optimization, ignoring the reality of hardware-target heterogeneity (e.g., deploying on a system that supports symmetric per-tensor rather than asymmetric per-channel representations). Second, they assume pristine, clean, and class-balanced calibration streams, masking the fragility of unsupervised entropy objectives under realistic test-time adaptation shifts \cite{mummadi2021tta}. Our work is deeply inspired by critical audits in related subfields, such as the LoRA-SAM audit of weight flatness \cite{lorasam2025}, and the overfitting-optimizer paradox investigations \cite{polymerge2025}. We depart from the standard "state-of-the-art" chasing narrative and instead construct a comprehensive, multi-axial deconstruction framework to rigorously map the true generalization and robustness boundaries of this emerging paradigm.
